% =====================================================================
% Bagian 6 — Responsible AI dan Keterbatasan  (target 1,25 halaman)
% Enam kategori rubrik disajikan sebagai tabel padat, bukan tujuh paragraf.
% Kolom "Status" WAJIB dipertahankan — ia yang memisahkan mitigasi yang
% sudah berjalan dari yang baru rancangan.
% =====================================================================

\section{Responsible AI dan Keterbatasan}
\label{sec:keterbatasan}

\subsection{Risiko dan mitigasi}

Sistem ini mengeluarkan penilaian yang memicu pengerahan tim ke lapangan, pada
lanskap tempat kepemilikan lahan, tanggung jawab hukum, dan kepentingan
masyarakat bersinggungan.

Enam kategori risiko beserta mitigasinya dirinci pada Tabel~\ref{tab:risiko}
(Lampiran~\ref{app:risiko}), dengan kolom status yang memisahkan mitigasi yang
\textbf{sudah berjalan} pada purwarupa dari yang \textbf{baru rancangan}; tanpa
pembedaan itu, daftar mitigasi akan terbaca sebagai perlindungan yang sudah ada.
Yang sudah berjalan: kunci API tersimpan di sisi server dan tidak pernah masuk
log, galat tidak bocor ke layar operator, fusi menekan alarm palsu termal ke nol,
dan keputusan akhir tetap di tangan operator. Yang masih rancangan: pengaburan
wajah dan pelat kendaraan, autentikasi papan pantau, kontrol akses berbasis
peran, log akses, serta pembatasan formal terhadap penggunaan ulang data untuk
pengawasan masyarakat.

Tiga catatan memerlukan penegasan. \textbf{Halusinasi} tidak berlaku karena model
bersifat diskriminatif; risiko analognya, keyakinan berlebih di luar distribusi
latih, memberi indikasi awal yang menenangkan: pada jalur termal-saja, keyakinan
rata-rata saat model keliru adalah $0{,}691$ dibanding $0{,}994$ saat benar.
\textbf{Keselamatan} didukung bukti langsung: \emph{recall} deteksi kejadian api
bernilai $1{,}000$ pada seluruh partisi dan level degradasi, dan seluruh
kesalahan berupa alarm palsu yang mekanismenya teridentifikasi
(Subbagian~4.9), risiko yang nyata untuk gambut, karena permukaan gambut kering
terbuka memanas ekstrem pada siang hari. \textbf{Generalisasi \emph{zero-shot}
92\% ke FLAME~3 bukan bukti transfer ke gambut tropis}: kedua dataset berada di
Amerika Serikat dan beriklim sedang.

Beberapa keputusan antarmuka ditujukan khusus agar sistem tidak tampak lebih
matang daripada kenyataannya, dan seluruhnya sudah berjalan: sensor darat
berlabel "simulasi" di setiap titik kemunculannya; blok latensi perangkat tepi
menampilkan "---" alih-alih angka rekaan; keadaan data yang belum dianalisis
tampil sebagai label eksplisit, bukan nilai kosong yang menyerupai galat;
catatan provenans menyatakan bingkai citra berasal dari FLAME~2 (Arizona) dengan
koordinat penempatan simulasi; dan penanda sumber data membedakan data karangan,
keluaran model, serta layanan hidup, ditulis otomatis oleh skrip inferensi
sehingga tidak dapat tertinggal basi. Sistem memberi prioritas, bukan
keputusan, operator menetapkan tindak lanjut melalui tiga pilihan yang terekam
pada jejak keputusan.

\subsection{Keterbatasan}

Ancaman terhadap validitas eksperimen diuraikan pada Subbagian~4.11; berikut
keterbatasan pada tingkat sistem dan penerapan. \textbf{Domain:} tidak ada satu
pun bingkai gambut tropis dalam seluruh pelatihan maupun evaluasi; sistem
dirancang untuk rezim membara dan divalidasi pada proksi terukur, dan belum
tervalidasi untuk gambut. \textbf{Variabel lingkungan:} kecepatan angin, suhu,
kelembapan, dan jenis vegetasi tidak menjadi masukan, keterbatasan yang
diidentifikasi narasumber ahli, sehingga keluaran berlaku pada asumsi kondisi
yang tidak diverifikasi. \textbf{Pemicu satelit:} data bersifat \emph{near
real-time} beresolusi 375~m, terhalang tutupan awan, menandai anomali termal
alih-alih kebakaran, dan belum divalidasi silang dengan keluaran model
(Lampiran~\ref{app:firms}). \textbf{Komponen yang belum nyata:} sensor darat
simulasi, rute patroli contoh, inferensi \emph{batch} luring, dan keputusan
operator tersimpan di memori proses. \textbf{Validasi pengguna:} tiga partisipan
mahasiswa alih-alih operator posko dan tanpa tekanan waktu, satu narasumber
ahli, serta evaluasi heuristik oleh pembangun produk sendiri. \textbf{Pengukuran
tepi:} tidak ada perangkat tepi fisik yang diuji, dan dampak kuantisasi terhadap
akurasi tidak dapat ditetapkan.

Kami memilih menyatakan keterbatasan ini secara rinci alih-alih meringkasnya.
Sistem yang menentukan ke mana tim pemadam dikerahkan tidak layak dinilai dari
klaim yang tidak dapat ditelusuri.
